\section{Kết luận và hướng phát triển}
\label{sec:ket-luan}

\subsection{Tóm tắt kết quả}
\label{sec:tom-tat-ket-qua}

\TODO{Tóm tắt lại phát hiện chính từ Mục~\ref{sec:ket-qua-chinh} và~\ref{sec:thao-luan}:
họ phương pháp nào giữ hiệu năng tốt nhất khi tỷ lệ thiếu dữ liệu tăng, và mức độ cải
thiện cụ thể so với các baseline ở Mục~\ref{sec:baseline}.}

Phần thực nghiệm SAITS đã hoàn tất 30 lượt huấn luyện trên 10 điều kiện thiếu
dữ liệu và ba seed. Macro F1 đạt $0.9224 \pm 0.0082$ khi không ẩn thêm; tại
mức thiếu 50\%, point, subsequence và block lần lượt đạt $0.8530 \pm 0.0071$,
$0.8030 \pm 0.0614$ và $0.8943 \pm 0.0061$. Trong cấu hình đã chạy, block
50\% giữ Macro F1 cao hơn hai kiểu thiếu còn lại ở cùng mức, trong khi
subsequence 50\% có độ biến thiên giữa seed lớn nhất. Đây là kết luận trong
phạm vi SAITS; việc xếp hạng giữa các họ mô hình cần kết quả đối chiếu và kiểm
định thống kê tương ứng.

\subsection{Hạn chế}
\label{sec:han-che}

\begin{itemize}
  \item \textbf{Nhãn từ API:} như đã nêu ở Mục~\ref{sec:han-che-nhan}, nhãn Rain/Snow
        được suy ra từ Dark Sky API, không phải quan trắc thực địa trực tiếp, nên mọi kết
        luận về độ chính xác tuyệt đối cần được hiểu là so với nhãn tham chiếu này.
  \item \textbf{Một trạm quan trắc:} toàn bộ benchmark chỉ dựa trên dữ liệu của một trạm
        (Szeged), nên khả năng tổng quát hóa sang vùng khí hậu khác chưa được kiểm chứng.
  \item \textbf{Mất dữ liệu là giả lập:} mặt nạ thiếu dữ liệu (Mục~\ref{sec:sinh-mat-na})
        được sinh tổng hợp, có thể không phản ánh đúng phân bố mất dữ liệu thực tế của
        một cảm biến cụ thể (ví dụ tương quan giữa các biến khi cùng một trạm gặp sự cố
        nguồn điện).
\end{itemize}

\subsection{Hướng phát triển}
\label{sec:huong-phat-trien}

Riêng SAITS, bước tiếp theo là ablation chuẩn hóa z-score trên phần quan sát
của train, giữ nguyên split và mask để đánh giá ảnh hưởng của thang đo đầu vào.
Cần kiểm tra độ ổn định của các điều kiện có biến thiên lớn, đặc biệt
subsequence 50\%, và bổ sung đánh giá imputation riêng nếu mở rộng mục tiêu
từ phân loại sang điền số đo thời tiết. Các thay đổi cấu hình cần được chọn
bằng validation, không sử dụng test để điều chỉnh siêu tham số.

\begin{itemize}
  \item \textbf{Neural CDE:} khai thác mô hình continuous-time (Neural Controlled
        Differential Equations) để xử lý tự nhiên dữ liệu lấy mẫu không đều, thay vì phải
        rời rạc hóa thành cửa sổ cố định 24 giờ.
  \item \textbf{Dữ liệu nhiều trạm:} mở rộng benchmark sang nhiều trạm quan trắc ở các
        vùng khí hậu khác nhau để kiểm chứng khả năng tổng quát hóa của kết luận.
  \item \textbf{Triển khai trên Edge:} đánh giá khả năng chạy mô hình tốt nhất (cân bằng
        giữa hiệu năng và chi phí tính toán, xem Mục~\ref{sec:chi-phi-tinh-toan}) trên
        thiết bị biên gắn tại trạm quan trắc thực tế.
\end{itemize}
